% !TeX root = ../../Book1.tex
% 纲要标题：### 5.3 小阶群的分析
% 纲要位置：计划纲要.md，第 261 行。

\section{小阶群的分析}
\label{b1:sec:0503}

在小阶群中，Sylow 子群个数的算术限制往往很强，但从个数走到群的乘法仍有距离。两素数乘积阶群提供了一个完整例子：唯一的 Sylow \(q\)-子群给出正规部分，另一个子群的共轭作用决定乘法。对于更多素因子的情形，有时只能排除单性。阅读这些例子时，应区分“必要条件”“已构造的群”和“全部同构类型”这三种结论。

% 本节正文按下列原纲要要点撰写。

% 纲要内容（仅作写作计划，尚非教材正文）：
%
% * pq 阶群及适当条件下的分类
% * 平方自由阶群的若干结论
% * 元素计数与单群排除
%

\subsection{两素数乘积阶群的分类}
\label{b1:subsec:0503:01}

本小节设 \(p<q\) 是不同素数。先由 Sylow 计数找到一个正规 \(q\) 阶子群，再考察 \(p\) 阶元素对它的共轭；这会把分类问题化为一个模 \(q\) 的参数。

\ProseParagraphBreak
分类证明的最后一步要比较不同参数，届时只需使用下面这个初等根数界。这里把多项式看作模 \(q\) 系数的形式表达式，证明只用显式恒等式和非零剩余类可逆，是一个局部自足的工具，不需要一般多项式环理论。

\begin{lemma}[模素数的多项式根数界]\label{b1:lem:polynomial-root-bound-mod-prime}
设 \(q\) 为素数，\(d\geq0\) 为整数，并设
\[
f(X)=c_0+c_1X+\cdots+c_dX^d,
\qquad c_i\in\mathbb Z/q\mathbb Z\ (0\leq i\leq d),\quad c_d\neq0.
\]
这里的多项式指由这些系数给出的形式表达式，条件 \(c_d\neq0\) 表明其次数为 \(d\)；“非零”指系数不全为零，而不是说它在剩余类集合上给出的函数不是零函数。那么，模 \(q\) 剩余类中至多有 \(d\) 个元素 \(c\) 满足 \(f(c)=0\)。
\end{lemma}
\begin{proof}
对 \(d\) 归纳。若 \(d=0\)，则 \(f=c_0\neq0\)，所以没有根。以下设 \(d\geq1\)。如果 \(f\) 没有根，结论已经成立；否则取一个根 \(c\)。由 \(f(c)=0\) 和逐项恒等式
\[
X^r-c^r=(X-c)(X^{r-1}+X^{r-2}c+\cdots+c^{r-1})
\]
可将 \(f\) 写成
\[
f(X)=f(X)-f(c)=(X-c)g(X).
\]
多项式 \(g\) 的首项为 \(c_dX^{d-1}\)，故 \(\deg g=d-1\)。若 \(e\neq c\) 也是 \(f\) 的根，则 \((e-c)g(e)=0\)。因为 \(q\) 是素数，非零剩余类 \(e-c\) 可逆，所以 \(g(e)=0\)。归纳假设说明这样的 \(e\) 至多有 \(d-1\) 个；再计入 \(c\)，\(f\) 至多有 \(d\) 个根。
\end{proof}

\begin{theorem}[两素数乘积阶群的分类]\label{b1:thm:groups-order-pq}
设 \(p<q\) 为素数。
\begin{enumerate}
\item 若 \(p\nmid(q-1)\)，则每个 \(pq\) 阶群都同构于循环群 \(C_{pq}\)。
\item 若 \(p\mid(q-1)\)，则在同构意义下恰有两个 \(pq\) 阶群：循环群 \(C_{pq}\) 和一个非交换群。后者可用剩余类对 \((i,j)\in(\mathbb Z/q\mathbb Z)\times(\mathbb Z/p\mathbb Z)\) 表示，乘法为
\[
(i,j)(k,\ell)=(i+t^jk,j+\ell),
\]
其中 \(t\) 是模 \(q\) 乘法群中一个阶为 \(p\) 的元素。
\end{enumerate}
\end{theorem}
\begin{proof}
设 \(|G|=pq\)。\cref{b1:thm:sylow-count}给出 \(n_q\mid p\) 且 \(n_q\equiv1\pmod q\)，因为 \(p<q\)，只能有 \(n_q=1\)。令 \(Q\) 是这个唯一的 Sylow \(q\)-子群，则同一定理保证 \(Q\trianglelefteq G\)。再由\cref{b1:thm:sylow-existence}取一个 \(p\) 阶子群 \(P\)。\cref{b1:cor:finite-order-divides}说明素数阶群循环，故可写 \(Q=\langle a\rangle\)、\(P=\langle b\rangle\)。

由\cref{b1:thm:lagrange}，交 \(Q\cap P\) 的阶同时整除 \(p,q\)，只能为 \(1\)。于是若 \(a^ib^j=a^kb^\ell\)，则 \(a^{i-k}=b^{\ell-j}\in Q\cap P\)，所以 \(i\equiv k\pmod q\)、\(j\equiv\ell\pmod p\)。取 \(0\leq i<q\)、\(0\leq j<p\)，便得到 \(pq\) 个不同的乘积，已遍历 \(G\)。因此每个元素有唯一的标准表达式 \(a^ib^j\)。

由于 \(Q\) 正规，对每个 \(u\in P\)，共轭 \(c_u(x)=uxu^{-1}\) 都是 \(Q\) 的自同构，逆映射为 \(c_{u^{-1}}\)。对 \(u,v\in P\) 又有 \(c_{uv}=c_u\circ c_v\)，得到同态
\[
P\longrightarrow\operatorname{Aut}(Q),\qquad u\longmapsto c_u.
\]
\(P\) 由 \(b\) 生成，\(Q\) 由 \(a\) 生成，所以整个共轭作用由 \(b\) 把 \(a\) 送到何处决定。写 \(bab^{-1}=a^t\)，其中 \(t\) 是模 \(q\) 的剩余类。自同构必须把生成元送到生成元；若 \(t=0\)，则 \(a^t=1\)，不可能是非单位元素 \(a\) 的像。因此 \(t\ne0\)，而 \(q\) 为素数也保证每个这样的 \(a^t\) 都生成 \(Q\)。

反复共轭得到 \(b^jab^{-j}=a^{t^j}\)。取 \(j=p\)，由 \(b^p=1\) 可得
\[
a=b^pab^{-p}=a^{t^p},\qquad a^{t^p-1}=1.
\]
取 \(t\) 的一个整数代表 \(T\)。因为 \(a\) 的阶为 \(q\)，从 \(a^{T^p-1}=1\) 得到 \(q\mid T^p-1\)，所以剩余类 \(t\) 满足 \(t^p=1\)，也即 \(t^p\equiv1\pmod q\)。更一般地，
\[
b^ja^kb^{-j}=(b^jab^{-j})^k=a^{t^jk}.
\]
乘积中的 \(b^j\) 和 \(a^k\) 通常不能交换；用这个共轭公式把乘积重新写成标准表达式，便有
\[
(a^ib^j)(a^kb^\ell)
=a^i(b^ja^kb^{-j})b^{j+\ell}
=a^{\,i+t^jk}b^{\,j+\ell}.
\]
因此 \(G\) 的全部乘法由 \(t\) 确定。若 \(t=1\)，则 \(a,b\) 交换，因而对正整数 \(d\) 有 \((ab)^d=a^db^d\)。若 \((ab)^d=1\)，则 \(a^d=b^{-d}\in Q\cap P=\{1\}\)，所以 \(q\mid d\) 且 \(p\mid d\)。又有 \((ab)^{pq}=a^{pq}b^{pq}=1\)，故 \(ab\) 的阶为 \(pq\)，\(G\) 循环。

由\cref{b1:thm:cyclic-automorphisms}，\(Q\) 的自同构群是阶为 \(q-1\) 的非零模 \(q\) 剩余类乘法群。若 \(t\neq1\) 且 \(t^p=1\)，则 \(t\) 阶为 \(p\)，对其生成的 \(p\) 阶子群使用\cref{b1:thm:lagrange}，得到 \(p\mid(q-1)\)。这证明第一项。

现在反过来构造群：把标准表达式中的指数对 \((i,j)\) 作为元素，将刚才算出的指数更新公式作为乘法。给定满足 \(t^p=1\) 的非零模 \(q\) 剩余类 \(t\)，在剩余类对
\[
(\mathbb Z/q\mathbb Z)\times(\mathbb Z/p\mathbb Z)
\]
上使用定理中的乘法，并将所得模型记为 \(G_t\)。第一坐标按模 \(q\) 运算，第二坐标按模 \(p\) 运算；因为 \(t^p=1\)，\(t^j\) 不依赖 \(j\) 的整数代表。因此运算良定义。对三个因子 \((i,j),(k,\ell),(r,s)\)，结合律的两边均为
\[
(i+t^jk+t^{j+\ell}r,\ j+\ell+s),
\]
单位是 \((0,0)\)，\((i,j)\) 的逆为 \((-t^{-j}i,-j)\)，所以 \(G_t\) 是一个 \(pq\) 阶群。对于前面任取的原群 \(G\) 及其共轭参数 \(t\)，定义
\[
\theta:G_t\longrightarrow G,
\qquad (i,j)\longmapsto a^ib^j.
\]
标准表达式的唯一性说明 \(\theta\) 是双射，而此前算出的乘法公式给出
\[
\theta\bigl((i,j)(k,\ell)\bigr)
=a^{\,i+t^jk}b^{\,j+\ell}
=(a^ib^j)(a^kb^\ell),
\]
所以 \(\theta\) 是群同构。由此，每个 \(pq\) 阶群都同构于由其参数得到的模型。

若 \(p\mid(q-1)\)，则可在阶为 \(q-1\) 的非零模 \(q\) 剩余类乘法群中应用\cref{b1:thm:cauchy}，取一个阶为 \(p\) 的元素 \(t\)。在 \(G_t\) 中取 \(a=(1,0)\)、\(b=(0,1)\)，便有 \(bab^{-1}=a^t\neq a\)，所以 \(G_t\) 非交换。这证明了非交换类型的存在性。

还需证明不同的非平凡参数不会产生新的同构类型。为此，先证明非零模 \(q\) 剩余类的乘法群中的 \(p\) 阶子群唯一，使不同参数只相差这个子群中生成元的选择。固定一个阶为 \(p\) 的元素 \(t\)，子群 \(\langle t\rangle\) 的 \(p\) 个元素 \(1,t,\ldots,t^{p-1}\) 互不相同，并且都是形式多项式 \(X^p-1\) 的根。由\cref{b1:lem:polynomial-root-bound-mod-prime}，这个 \(p\) 次多项式不再有其他根。因此所有阶为 \(p\) 的元素都在 \(\langle t\rangle\) 中，任意 \(p\) 阶子群也只能等于 \(\langle t\rangle\)。特别地，对另一个阶为 \(p\) 的参数 \(t'\)，必有 \(t'=t^u\)，其中 \(u\not\equiv0\pmod p\)。定义
\[
F_u:G_{t'}\longrightarrow G_t,
\qquad (i,j)\longmapsto(i,uj).
\]
由于 \(u\) 在模 \(p\) 意义下非零，第二坐标上的乘法可逆，故 \(F_u\) 是双射。若分别以 \(*_{t'}\) 和 \(*_t\) 表示两个模型的乘法，则
\[
\begin{aligned}
F_u\bigl((i,j)*_{t'}(k,\ell)\bigr)
  &=\bigl(i+t'^{\,j}k,\,u(j+\ell)\bigr)\\
  &=\bigl(i+t^{uj}k,\,uj+u\ell\bigr)\\
  &=F_u(i,j)*_tF_u(k,\ell).
\end{aligned}
\]
因此 \(F_u\) 是群同构，所有非交换情形合并为一个同构类。这里唯一的是非交换群的同构类型，不是参数 \(t\)、生成元 \(a,b\) 或具体同构映射；例如将 \(b\) 换成 \(b^u\)，共轭参数就从 \(t\) 变为 \(t^u\)。
\end{proof}

这里用的是模素数剩余类的初等运算和根因子恒等式，并未假设后文的有限域乘法群循环性。上述乘法将在\secref{b1:sec:0602}中作为半直积的特例重新解释。例如 \(21=3\cdot7\)，可取 \(t=2\)，因为 \(2^3\equiv1\pmod7\) 而 \(2\not\equiv1\pmod7\)；于是非交换 \(21\) 阶群存在，且在同构意义下唯一。

\subsection{平方自由阶群的若干结论}
\label{b1:subsec:0503:02}

设 \(n\) 为正整数。如果不存在素数 \(r\) 使 \(r^2\mid n\)，那么称 \(n\) 为\term[pingfang-ziyou]{平方自由}[squarefree]；等价地，\(n\) 的素因子分解中每个素因子的指数都是一。若 \(|G|\) 平方自由，则对每个整除 \(|G|\) 的素数 \(p\)，Sylow \(p\)-子群都有素数阶，因而循环。不过，循环 Sylow 子群不保证整个群循环，\(S_3\) 就是最小的反例。

\begin{proposition}\label{b1:prop:squarefree-normal-sylow}
设有限群 \(G\) 的阶平方自由。如果它的每个 Sylow 子群都正规，则 \(G\) 循环。
\end{proposition}
\begin{proof}
平凡群的情形显然，以下设 \(|G|>1\)。先证明不同 Sylow 子群逐元素交换。若 \(P,Q\) 是对应不同素数的正规 Sylow 子群，取 \(x\in P,y\in Q\)。由 \(Q\) 正规，\(xyx^{-1}\in Q\)，再乘上 \(y^{-1}\in Q\)，得到
\[
xyx^{-1}y^{-1}=(xyx^{-1})y^{-1}\in Q.
\]
另一方面，由 \(P\) 正规，\(yx^{-1}y^{-1}\in P\)，再乘上 \(x\in P\)，得到
\[
xyx^{-1}y^{-1}=x(yx^{-1}y^{-1})\in P.
\]
因此这个交换子属于 \(P\cap Q\)。由\cref{b1:thm:lagrange}，交的阶同时整除这两个互素的子群阶，所以 \(P\cap Q=\{1\}\)，交换子只能为单位元，故 \(xy=yx\)。

把各 Sylow 子群记为 \(P_i=\langle a_i\rangle\)，其阶为不同素数 \(p_i\)；素数阶群循环的依据是\cref{b1:cor:finite-order-divides}。下面说明这些子群的乘积等于 \(G\)。由于各因子逐元素交换，有限乘积 \(P_1\cdots P_j\) 对乘法和取逆封闭，是子群。其阶可逐步计算：若 \(|P_1\cdots P_{j-1}|=p_1\cdots p_{j-1}\)，则它与 \(P_j\) 的交阶同时整除 \(p_1\cdots p_{j-1}\) 和 \(p_j\)，故交平凡。若 \(uh=u'h'\)，其中 \(u,u'\in P_1\cdots P_{j-1}\)、\(h,h'\in P_j\)，则 \(u'^{-1}u=h'h^{-1}\) 属于该交，所以每个乘积都有唯一的两个因子，阶便乘上 \(p_j\)。由此归纳得到全部乘积子群的阶为 \(\prod_i p_i=|G|\)，因而它等于 \(G\)。

最后取 \(a=\prod_i a_i\)。若 \(a^d=1\)，将第 \(i\) 个因子分离，则 \(a_i^d\) 同时属于阶为 \(p_i\) 的子群和由其他 Sylow 子群组成的、阶与 \(p_i\) 互素的子群，故 \(p_i\mid d\)。各 \(p_i\) 互素，得到 \(|G|=\prod_i p_i\mid d\)。另一方面，由\cref{b1:thm:lagrange}，\(a\) 的阶整除 \(|G|\)，所以其阶恰为 \(|G|\)，从而生成 \(G\)。
\end{proof}

上述证明还给出，各正规 Sylow 子群逐元素交换，乘积是整个群，而且每个因子与其余因子的乘积相交平凡。因此，每个群元素都有唯一的逐因子表达式。\secref{b1:sec:0601}将把这种结构表述为内直积。

\ProseParagraphBreak
这个命题保留了“全部正规”这一实质条件，并未给出任意平方自由阶群的完整分类；当素因子增多时，正规性和不同子群间的共轭作用还需要进一步分析。

\subsection{元素计数与单群排除}
\label{b1:subsec:0503:03}

对于阶为合数的单群，任何 Sylow 子群只要非平凡且为真子群，就不能唯一。这往往能迅速排除许多群阶。

\ProseParagraphBreak
若 \(|G|=28\)，则\cref{b1:thm:sylow-count}给出 \(n_7\mid4\)、\(n_7\equiv1\pmod7\)，所以 \(n_7=1\)，得到正规七阶子群，故 \(G\) 不是单群。若单个 Sylow 子群个数尚不能判定，就须同时考虑不同素数对应的元素。

\ProseParagraphBreak
只有确定子群间的交以后，才能判断各子群的元素数是否可以直接相加。记 \(c_p\) 为 \(G\) 的全部 \(p\) 阶子群的个数。每个 \(p\) 阶元素生成唯一的 \(p\) 阶子群，而不同的这类子群除单位元外互不相交，因此 \(G\) 中 \(p\) 阶元素的总数总是
\[
c_p(p-1).
\]
下面的 \(30\) 阶例子需同时计算三阶元素和五阶元素。

\begin{proposition}\label{b1:prop:order-thirty-not-simple}
每个 \(30\) 阶群都含有正规的三阶或五阶子群，因而不是单群。
\end{proposition}
\begin{proof}
\cref{b1:thm:sylow-count}给出 \(n_5\mid6\)、\(n_5\equiv1\pmod5\)，以及 \(n_3\mid10\)、\(n_3\equiv1\pmod3\)。若这两类 Sylow 子群都不唯一，则只能为 \(n_5=6\)、\(n_3=10\)。由于 \(30\) 中的素因子 \(3,5\) 都只出现一次，这些 Sylow 子群分别是三阶和五阶子群。按照上面的计数，每一类除单位外互不相交，故分别有 \(6(5-1)=24\) 个五阶元素和 \(10(3-1)=20\) 个三阶元素。两种元素阶不同，彼此也不重叠；连同单位超过 \(30\) 个，矛盾。故至少有 \(n_3=1\) 或 \(n_5=1\)，由\cref{b1:thm:sylow-count}，相应子群正规。它的阶为 \(3\) 或 \(5\)，是非平凡真子群，故群不是单群。
\end{proof}

这个例子可以直接用 \(n_3,n_5\) 计数，是因为三阶、五阶子群恰为相应的 Sylow 子群；一般情形不能无条件把 \(c_p\) 换成 \(n_p\)。当 \(p\bigm\vert|G|\) 且 \(p^2\mathrel{\vcenter{\hbox{\scalebox{1.25}{\(\nmid\)}}}}|G|\) 时，两类子群确实相同，因而 \(c_p=n_p\)。在其他情形中，必须另行证明 \(c_p=n_p\)，才能作这种替换。例如循环群 \(C_{p^2}\) 中二者都等于一，尽管 \(p^2\bigm\vert|C_{p^2}|\)。

\ProseParagraphBreak
例如 \(V_4\) 只有一个 Sylow \(2\)-子群，就是它自身，故 \(n_2=1\)；它却有三个二阶子群和三个二阶元素，不能用 \(n_2(2-1)=1\) 计数。还应区分“阶恰为 \(p\) 的元素”与“阶为某个 \(p^r\)（\(r\geq1\)）的元素”：若 Sylow 子群阶为 \(p^a\) 且 \(a>1\)，同一子群中可能出现不同的非平凡 \(p\) 幂阶。每个 \(p\) 幂阶元素都在某个 Sylow 子群中，这是将\cref{b1:thm:sylow-conjugacy}用于它生成的循环子群所得的结论；但不同 Sylow 子群可能共享非单位元素，故 \(n_p(p^a-1)\) 也未必等于全群非单位 \(p\) 幂阶元素的总数。下一节的置换作用提供一种不需要假设交平凡的排除方法。
